% MATERI INSTRUKSIONAL INDEPENDEN --- BUKAN TEKS SUMBER WEN-WEI LI.
% stable-unit-id: o014.mastery.derived-spectral.002
% provenance: OpenAI Codex gpt-5.6-sol, Ultra
% license: CC BY 4.0; kompatibel dengan lisensi terjemahan sumber.

\section*{Jembatan Penguasaan 002: Funktor Turunan dan Barisan Spektral}
\addcontentsline{toc}{section}{Jembatan Penguasaan 002: Funktor Turunan dan Barisan Spektral}
\hypertarget{o014.mastery.derived-spectral.002}{ }
\label{mastery:derived-spectral-002}

\noindent\textbf{Status dan atribusi.}
Unit ini merupakan materi instruksional independen yang disusun untuk
mendampingi buku ini. Unit ini bukan bagian dari teks sumber Wen-Wei Li dan
tidak boleh dinisbatkan kepadanya. Provenans penyusunan:
\texttt{OpenAI Codex gpt-5.6-sol, Ultra}. Materi ini tersedia berdasarkan
Creative Commons Attribution 4.0 International
(CC BY 4.0; \url{https://creativecommons.org/licenses/by/4.0/}), sehingga
kompatibel dengan lisensi teks sumber.

\subsection*{Peta prasyarat}

Rantai prasyarat yang paling singkat adalah sebagai berikut.
\begin{compactitem}
	\item Kategori abelian, kompleks, dan kohomologi:
	\S\ref{sec:Abel-cat-def}, \S\ref{sec:cohomology}, dan
	Definisi~\ref{def:cohomology}.
	\item Objek injektif/projektif serta keberadaan cukup banyak objek semacam itu:
	\S\ref{sec:inj-proj}, Definisi~\ref{def:injective-projective-obj}, dan
	Definisi~\ref{def:enough-injectives}.
	\item Resolusi dan keunikan hingga homotopi:
	\S\ref{sec:resolutions}, Definisi~\ref{def:resolutions}, serta
	Teorema~\ref{prop:resolution-homotopy}.
	\item Funktor turunan klasik, barisan eksak panjang, dan pergeseran
	dimensi: \S\ref{sec:derived-primer},
	Teorema~\ref{prop:long-exact-sequence-primer}, dan
	Proposisi~\ref{prop:dimension-shifting}.
	\item Barisan spektral, tepi, dan komposisi funktor:
	\S\ref{sec:spectral-sequence}, Catatan~\ref{rem:edge-computing}, dan
	Teorema~\ref{prop:Grothendieck-ss}.
	\item Kohomologi grup siklik dan struktur perkalian:
	Definisi~\ref{def:group-coh-ho},
	Proposisi~\ref{prop:group-cohomology-cyclic},
	\S\ref{sec:grading-filtration}, serta
	\S\ref{sec:multiplicative-SS}.
\end{compactitem}
Alur konseptualnya ialah
\[
	\text{ganti dengan resolusi}
	\Longrightarrow \text{ambil kohomologi}
	\Longrightarrow \text{susun filtrasi total}
	\Longrightarrow \text{baca halaman dan tepi}.
\]

\subsection*{Metode kerja: tiga buku besar}

Agar perhitungan tetap tertib, informasinya dapat diatur dalam tiga ``buku
besar'' yang tidak boleh dicampur.
\begin{enumerate}
	\item \emph{Buku resolusi} mencatat objek yang diganti dengan resolusi, arah resolusi,
	dan alasan bahwa suku-sukunya injektif, projektif, atau asiklik bagi funktor
	yang akan diterapkan.
	\item \emph{Buku derajat} mencatat indeks kohomologis, total derajat
	$n=p+q$, arah diferensial
	$d_r:E_r^{p,q}\to E_r^{p+r,q-r+1}$, dan semua suku yang lenyap.
	\item \emph{Buku filtrasi} mencatat objek abutmen $H^n$, filtrasi
	$\mathrm{F}^pH^n$, identifikasi
	$\gr^pH^n\simeq E_\infty^{p,n-p}$, serta kompatibilitas perkalian.
\end{enumerate}
Kesalahan yang paling umum ialah menyamakan $E_2^{p,q}$ langsung dengan
subobjek abutmen. Yang tersedia secara kanonik ialah suku bergradasi
$E_\infty^{p,q}$; untuk kembali ke abutmen, masalah ekstensi filtrasi masih
harus diselesaikan.

Sebagai contoh kerja, ambil barisan spektral kohomologis kuadran pertama
\[
	E_2^{p,q}\Rightarrow H^{p+q}.
\]
Dalam total derajat $1$, tidak ada diferensial yang menyentuh
$E_2^{1,0}$, sedangkan satu-satunya diferensial yang dapat keluar dari
$E_2^{0,1}$ ialah
$d_2:E_2^{0,1}\to E_2^{2,0}$. Karena itu,
\[
	E_\infty^{1,0}=E_2^{1,0},\qquad
	E_\infty^{0,1}=\Ker(d_2:E_2^{0,1}\to E_2^{2,0}).
\]
Filtrasi $H^1$ mempunyai subobjek
$\mathrm{F}^1H^1\simeq E_\infty^{1,0}$ dan hasil bagi
$H^1/\mathrm{F}^1H^1\simeq E_\infty^{0,1}$. Pada total derajat $2$,
\[
	E_\infty^{2,0}\simeq
	E_2^{2,0}/\Image(d_2:E_2^{0,1}\to E_2^{2,0})
	\hookrightarrow H^2.
\]
Ketiga catatan ini langsung menghasilkan barisan lima suku
\[
	0\to E_2^{1,0}\to H^1\to E_2^{0,1}
	\xrightarrow{d_2}E_2^{2,0}\to H^2.
\]
Metodenya selalu sama: tandai daerah yang lenyap, tentukan diferensial yang masih
mungkin, hitung $E_\infty$, baru kemudian gunakan filtrasi abutmen.

\subsection*{Latihan dan solusi lengkap}

% stable-exercise-id: o014.mastery.derived-spectral.002.ex01
\subsubsection*{Latihan 1 --- Perbandingan resolusi dan independensi pilihan}
\hypertarget{o014.mastery.derived-spectral.002.ex01}{ }
\label{mastery:derived-spectral-002-ex01}

Misalkan $\mathcal{A}$ adalah kategori abelian dengan cukup banyak objek injektif,
$F:\mathcal{A}\to\mathcal{B}$ adalah funktor aditif, dan
\[
	0\to X\to I^0\to I^1\to\cdots,
	\qquad
	0\to X\to J^0\to J^1\to\cdots
\]
dua resolusi injektif dari $X$.
\begin{enumerate}[(i)]
	\item Bangun pemetaan kompleks $u:I^\bullet\to J^\bullet$ dan
	$v:J^\bullet\to I^\bullet$ yang mengangkat $\identity_X$.
	\item Buktikan bahwa $vu\simeq\identity_I$ dan
	$uv\simeq\identity_J$ melalui suatu homotopi rantai.
	\item Simpulkan bahwa $\Hm^n(FI^\bullet)$ dan
	$\Hm^n(FJ^\bullet)$ isomorfik secara kanonik untuk setiap $n$.
\end{enumerate}

% stable-solution-id: o014.mastery.derived-spectral.002.sol01
\paragraph{Solusi.}
\hypertarget{o014.mastery.derived-spectral.002.sol01}{ }
Karena $X\hookrightarrow I^0$ monomorfisme dan $J^0$ injektif, morfisme
$X\to J^0$ memanjang menjadi $u^0:I^0\to J^0$. Andaikan
$u^0,\ldots,u^n$ telah dibangun dan memenuhi syarat pemetaan kompleks.
Morfisme $d_J^nu^n:I^n\to J^{n+1}$ bernilai nol pada
$\Ker(d_I^n)=\Image(d_I^{n-1})$, sehingga berfaktor melalui
$\Image(d_I^n)\hookrightarrow I^{n+1}$. Karena $J^{n+1}$ injektif,
faktorisasi itu memanjang menjadi $u^{n+1}:I^{n+1}\to J^{n+1}$ dengan
\[
	u^{n+1}d_I^n=d_J^nu^n.
\]
Induksi menghasilkan $u$. Dengan menukar $I$ dan $J$, diperoleh $v$.

Kita juga memerlukan fakta bahwa dua pengangkatan $r,s:I^\bullet\to
J^\bullet$ dari $\identity_X$ homotopik. Pada derajat nol, $r^0-s^0$
bernilai nol pada $X=\Ker(d_I^0)$, sehingga berfaktor melalui
$\Image(d_I^0)\hookrightarrow I^1$. Injektivitas $J^0$ memberikan
$h^1:I^1\to J^0$ dengan
$r^0-s^0=h^1d_I^0$. Secara induktif, setelah $h^1,\ldots,h^n$ dibangun,
morfisme
\[
	r^n-s^n-d_J^{n-1}h^n:I^n\longrightarrow J^n
\]
bernilai nol pada $\Ker(d_I^n)$. Ia berfaktor melalui
$\Image(d_I^n)\hookrightarrow I^{n+1}$ dan, karena $J^n$ injektif,
memanjang menjadi $h^{n+1}:I^{n+1}\to J^n$. Jadi
\[
	r-s=d_Jh+hd_I,
\]
yakni $r\simeq s$.

Komposisi $vu$ dan $\identity_I$ sama-sama mengangkat $\identity_X$, maka
keduanya homotopik; demikian pula $uv\simeq\identity_J$. Karena $F$ aditif,
$F$ membawa persamaan homotopi ke persamaan homotopi. Oleh sebab itu,
$\Hm^n(Fu)$ dan $\Hm^n(Fv)$ saling invers. Jika pemetaan perbandingan lain
dipilih, pemetaan itu homotopik dengan yang pertama dan menginduksi pemetaan kohomologi
yang sama. Jadi isomorfisme tersebut kanonik. Inilah isi konkret independensi
resolusi dalam Teorema~\ref{prop:resolution-homotopy}.

% stable-exercise-id: o014.mastery.derived-spectral.002.ex02
\subsubsection*{Latihan 2 --- Menghitung $\Ext$ dari resolusi pendek}
\hypertarget{o014.mastery.derived-spectral.002.ex02}{ }
\label{mastery:derived-spectral-002-ex02}

Misalkan $m\geq1$ dan $M$ adalah grup abelian. Gunakan resolusi projektif
\[
	0\longrightarrow\Z\xrightarrow{\times m}\Z
	\longrightarrow\Z/m\Z\longrightarrow0
\]
untuk menghitung $\Ext^n_{\Z}(\Z/m\Z,M)$ untuk semua $n\geq0$.
Jelaskan mengapa hasilnya tidak bergantung pada resolusi projektif yang
dipilih.

% stable-solution-id: o014.mastery.derived-spectral.002.sol02
\paragraph{Solusi.}
\hypertarget{o014.mastery.derived-spectral.002.sol02}{ }
Terapkan $\Hom_{\Z}(-,M)$ pada resolusi tersebut. Karena
$\Hom_{\Z}(\Z,M)\simeq M$, diperoleh kompleks kokrantai
\[
	0\longrightarrow
	\underbracket{M}_{\text{derajat }0}
	\xrightarrow{\times m}
	\underbracket{M}_{\text{derajat }1}
	\longrightarrow0.
\]
Jadi
\[
	\Ext^n_{\Z}(\Z/m\Z,M)\simeq
	\begin{cases}
		M[m]:=\{x\in M:mx=0\},&n=0,\\
		M/mM,&n=1,\\
		0,&n\geq2.
	\end{cases}
\]
Untuk $n=0$, identifikasi ini sama dengan
$\Hom_{\Z}(\Z/m\Z,M)\simeq M[m]$. Dua resolusi projektif dari
$\Z/m\Z$ dihubungkan oleh pemetaan perbandingan yang unik hingga homotopi,
yakni versi dual Latihan 1. Funktor $\Hom_{\Z}(-,M)$ memetakan homotopi
tersebut menjadi homotopi, sehingga kohomologinya tidak berubah. Karena itu, hasil di atas
tidak bergantung pada resolusi yang digunakan.

% stable-exercise-id: o014.mastery.derived-spectral.002.ex03
\subsubsection*{Latihan 3 --- Barisan eksak panjang funktor turunan}
\hypertarget{o014.mastery.derived-spectral.002.ex03}{ }
\label{mastery:derived-spectral-002-ex03}

Misalkan $F:\mathcal{A}\to\mathcal{B}$ adalah funktor aditif eksak kiri dan
$\mathcal{A}$ mempunyai cukup banyak objek injektif. Untuk barisan eksak
pendek
\[
	0\longrightarrow X\longrightarrow Y\longrightarrow Z\longrightarrow0,
\]
gunakan resolusi injektif yang serasi untuk membangun morfisme penghubung
\[
	\delta^n:\mathrm{R}^nF(Z)\longrightarrow\mathrm{R}^{n+1}F(X).
\]
Buktikan bahwa morfisme tersebut terdefinisi dengan baik, alami, dan menghasilkan
barisan eksak panjang.

% stable-solution-id: o014.mastery.derived-spectral.002.sol03
\paragraph{Solusi.}
\hypertarget{o014.mastery.derived-spectral.002.sol03}{ }
Versi injektif lema tapal kuda memberikan diagram resolusi yang baris-barisnya
eksak dan setiap derajatnya terbelah:
\[
\begin{tikzcd}
	0 \arrow[r] & X \arrow[r] \arrow[d] & Y \arrow[r] \arrow[d] & Z \arrow[r] \arrow[d] & 0 \\
	0 \arrow[r] & I_X^\bullet \arrow[r] & I_Y^\bullet \arrow[r] & I_Z^\bullet \arrow[r] & 0.
\end{tikzcd}
\]
Karena barisan bawah terbelah pada setiap derajat, setiap funktor aditif
mempertahankan keeksakannya. Jadi
\[
	0\to F(I_X^\bullet)\to F(I_Y^\bullet)
	\to F(I_Z^\bullet)\to0
\]
merupakan barisan eksak pendek kompleks.

Ambil kelas $[z]\in\Hm^n(FI_Z^\bullet)$ dengan wakil kokitar $z$.
Angkat $z$ menjadi $y\in F(I_Y^n)$. Karena citra $d(y)$ di
$F(I_Z^{n+1})$ sama dengan $d(z)=0$, terdapat $x\in F(I_X^{n+1})$ dengan
citra $x$ sama dengan $d(y)$. Karena pemetaan kompleks di sebelah kiri injektif dan
$d^2=0$ menunjukkan $d(x)=0$. Definisikan
\[
	\delta^n[z]:=[x].
\]
Jika $y$ diganti oleh pengangkatan lain, selisihnya berasal dari
$F(I_X^n)$ dan $x$ berubah sebesar suatu kokobatas. Jika $z$ diganti oleh
$z+d(w)$, angkat $w$ dan koreksi $y$ dengan diferensial pengangkatan itu;
kelas $[x]$ tetap sama. Jadi $\delta^n$ terdefinisi dengan baik.

Keeksakan diperiksa dengan pengejaran yang sama. Sebagai contoh,
$[y]\in\Hm^n(FI_Y)$ dipetakan ke nol di $\Hm^n(FI_Z)$ jika dan hanya jika,
setelah $y$ dikoreksi oleh suatu kokobatas, ia berasal dari kokitar di
$F(I_X^n)$. Selanjutnya, $[z]$ berada dalam kernel $\delta^n$ jika dan hanya
jika $x$ suatu kokobatas; mengoreksi pengangkatan $y$ dengan suatu prapeta dari
kokobatas tersebut menghasilkan kokitar yang memetakan ke $z$. Argumen
serupa di suku berikutnya menyelesaikan keeksakan, sehingga diperoleh
\[
	\cdots\to\mathrm{R}^nF(X)\to\mathrm{R}^nF(Y)
	\to\mathrm{R}^nF(Z)\xrightarrow{\delta^n}
	\mathrm{R}^{n+1}F(X)\to\cdots.
\]
Untuk morfisme di antara dua barisan eksak pendek, pemetaan perbandingan
resolusi membawa pilihan $z,y,x$ menjadi pilihan yang sesuai dalam diagram kedua.
Kedua urutan operasi menghasilkan kelas yang sama; jadi $\delta^n$ alami.
Konstruksi ini mewujudkan Teorema~\ref{prop:long-exact-sequence-primer}.

% stable-exercise-id: o014.mastery.derived-spectral.002.ex04
\subsubsection*{Latihan 4 --- Pergeseran dimensi}
\hypertarget{o014.mastery.derived-spectral.002.ex04}{ }
\label{mastery:derived-spectral-002-ex04}

Misalkan $F:\mathcal{A}\to\mathcal{B}$ adalah funktor eksak kiri dan
\[
	0\longrightarrow X\longrightarrow A\longrightarrow B\longrightarrow0
\]
eksak, dengan $A$ bersifat $F$-asiklik.
\begin{enumerate}[(i)]
	\item Buktikan isomorfisme alami
	\[
		\mathrm{R}^1F(X)\simeq\Coker[FA\to FB],
		\qquad
		\mathrm{R}^nF(X)\simeq\mathrm{R}^{n-1}F(B)\qquad(n\geq2).
	\]
	\item Untuk resolusi injektif
	$0\to X\to I^0\to I^1\to\cdots$, definisikan
	$C^0=X$ dan $C^{r+1}=\Coker[C^r\to I^r]$. Gunakan (i) berulang kali
	untuk memperoleh rumus kohomologis untuk $\mathrm{R}^nF(X)$.
	\item Jika $X$ mempunyai resolusi injektif dengan $I^r=0$ untuk $r>d$,
	buktikan $\mathrm{R}^nF(X)=0$ untuk $n>d$.
\end{enumerate}

% stable-solution-id: o014.mastery.derived-spectral.002.sol04
\paragraph{Solusi.}
\hypertarget{o014.mastery.derived-spectral.002.sol04}{ }
Barisan eksak panjang dari Latihan 3 memuat
\[
	FA\longrightarrow FB\longrightarrow\mathrm{R}^1F(X)
	\longrightarrow\underbracket{\mathrm{R}^1F(A)}_{0}
\]
dan, untuk $n\geq2$,
\[
	\underbracket{\mathrm{R}^{n-1}F(A)}_{0}
	\longrightarrow\mathrm{R}^{n-1}F(B)
	\longrightarrow\mathrm{R}^nF(X)
	\longrightarrow\underbracket{\mathrm{R}^nF(A)}_{0}.
\]
Keeksakan memberikan kedua isomorfisme pada (i), dan sifat alaminya berasal
dari sifat alami barisan eksak panjang.

Resolusi injektif terurai menjadi barisan eksak pendek
\[
	0\to C^r\to I^r\to C^{r+1}\to0,
	\qquad r\geq0.
\]
Karena $I^r$ injektif, ia $F$-asiklik. Maka, untuk $n\geq1$,
\[
	\mathrm{R}^nF(X)\simeq\mathrm{R}^{n-1}F(C^1)\simeq\cdots
	\simeq\mathrm{R}^1F(C^{n-1})
	\simeq\Coker[FI^{n-1}\to FC^n].
\]
Keeksakan kiri $F$ mengidentifikasi $FC^n$ dengan
$\Ker[FI^n\to FI^{n+1}]$. Oleh sebab itu,
\[
	\mathrm{R}^nF(X)\simeq
	\frac{\Ker[FI^n\to FI^{n+1}]}
	{\Image[FI^{n-1}\to FI^n]}
	=\Hm^n(FI^\bullet).
\]
Jika resolusi berhenti setelah $I^d$, kompleks $FI^\bullet$ tidak mempunyai
suku berderajat $>d$, sehingga kohomologinya nol pada derajat $n>d$.
Pernyataan ini juga mengikuti dengan menggeser dimensi sampai kosizigi
terakhir yang injektif. Ini merupakan bentuk komputasional dari
Proposisi~\ref{prop:dimension-shifting}.

% stable-exercise-id: o014.mastery.derived-spectral.002.ex05
\subsubsection*{Latihan 5 --- Menyiapkan barisan spektral Grothendieck}
\hypertarget{o014.mastery.derived-spectral.002.ex05}{ }
\label{mastery:derived-spectral-002-ex05}

Misalkan
\[
	\mathcal{A}\xrightarrow{F}\mathcal{B}\xrightarrow{G}\mathcal{C}
\]
merupakan funktor-funktor aditif eksak kiri. Andaikan $\mathcal{A}$ dan
$\mathcal{B}$ mempunyai cukup banyak objek injektif, serta $F$ memetakan
setiap objek injektif ke objek $G$-asiklik. Untuk $X\in\Obj(\mathcal{A})$,
buktikan adanya barisan spektral kuadran pertama
\[
	E_2^{p,q}=\mathrm{R}^pG(\mathrm{R}^qF(X))
	\Rightarrow\mathrm{R}^{p+q}(GF)(X).
\]
Jelaskan secara eksplisit peran hipotesis $G$-asiklik.

% stable-solution-id: o014.mastery.derived-spectral.002.sol05
\paragraph{Solusi.}
\hypertarget{o014.mastery.derived-spectral.002.sol05}{ }
Pilih resolusi injektif $0\to X\to I^\bullet$. Kompleks
$F(I^\bullet)$ berada di kuadran tak negatif. Pilih resolusi
Cartan--Eilenberg $F(I^\bullet)\to J^{\bullet,\bullet}$ dalam
$\mathcal{B}$ dan terapkan $G$ untuk memperoleh bikompleks kuadran pertama
$GJ^{p,q}$. Dua filtrasi kompleks total
$\tot(GJ)$ memberikan dua barisan spektral yang berkonvergensi ke abutmen yang sama.

Untuk filtrasi pertama, ambil kohomologi dalam arah yang meresolusi setiap
$F(I^p)$. Diperoleh
\[
	\Hm^q(GJ^{p,\bullet})
	\simeq\mathrm{R}^qG(F(I^p)).
\]
Hipotesis $G$-asiklik tepat mengatakan bahwa suku ini nol untuk $q>0$,
sedangkan untuk $q=0$ sukunya $GF(I^p)$. Jadi barisan spektral tersebut hanya
mempunyai satu baris tak nol dan berdegenerasi. Dengan demikian, kohomologi totalnya
ialah
\[
	\Hm^n(GF(I^\bullet))=\mathrm{R}^n(GF)(X).
\]

Untuk filtrasi kedua, sifat resolusi Cartan--Eilenberg menyatakan bahwa
kompleks yang diperoleh setelah mengambil kohomologi pada arah
$F(I^\bullet)$ merupakan resolusi injektif dari
\[
	\Hm^q(F(I^\bullet))=\mathrm{R}^qF(X).
\]
Karena itu, setelah menerapkan $G$ dan mengambil kohomologi, diperoleh
\[
	E_2^{p,q}=\mathrm{R}^pG(\mathrm{R}^qF(X)).
\]
Kedua filtrasi menghitung kohomologi total yang sama, yang menurut perhitungan
pertama adalah $\mathrm{R}^{p+q}(GF)(X)$. Karena bikompleks berada di kuadran
pertama, filtrasi pada setiap total derajat berhingga dan konvergensi berlaku.
Inilah konstruksi dalam Teorema~\ref{prop:Grothendieck-ss}; tanpa hipotesis
$G$-asiklik, barisan spektral pertama tidak harus runtuh dan abutmen tidak
dapat diidentifikasi dengan funktor turunan komposisi melalui argumen ini.

% stable-exercise-id: o014.mastery.derived-spectral.002.ex06
\subsubsection*{Latihan 6 --- Morfisme tepi dan lima suku Grothendieck}
\hypertarget{o014.mastery.derived-spectral.002.ex06}{ }
\label{mastery:derived-spectral-002-ex06}

Gunakan barisan spektral pada Latihan 5.
\begin{enumerate}[(i)]
	\item Identifikasi kedua morfisme tepi
	\[
		\mathrm{R}^nG(FX)\longrightarrow\mathrm{R}^n(GF)(X),
		\qquad
		\mathrm{R}^n(GF)(X)\longrightarrow G(\mathrm{R}^nF(X)).
	\]
	\item Turunkan barisan eksak lima suku
	\[
	\begin{aligned}
		0\to\mathrm{R}^1G(FX)&\to\mathrm{R}^1(GF)(X)
		\to G(\mathrm{R}^1F(X))\\
		&\xrightarrow{d_2}\mathrm{R}^2G(FX)
		\to\mathrm{R}^2(GF)(X).
	\end{aligned}
	\]
	\item Buktikan bahwa morfisme tepi pertama merupakan isomorfisme untuk
	semua $n$ jika $\mathrm{R}^qF(X)=0$ untuk $q>0$, dan bahwa morfisme tepi
	kedua merupakan isomorfisme untuk semua $n$ jika
	$\mathrm{R}^pG(\mathrm{R}^qF(X))=0$ untuk $p>0$.
\end{enumerate}

% stable-solution-id: o014.mastery.derived-spectral.002.sol06
\paragraph{Solusi.}
\hypertarget{o014.mastery.derived-spectral.002.sol06}{ }
Pada tepi bawah,
$E_2^{n,0}=\mathrm{R}^nG(FX)$ memetakan secara surjektif, melalui halaman
demi halaman, ke $E_\infty^{n,0}$. Suku terakhir ini adalah bagian
filtrasi terdalam $\mathrm{F}^n\mathrm{R}^n(GF)(X)$, sehingga inklusinya ke
abutmen memberi morfisme tepi pertama. Pada tepi kiri, hasil bagi
\[
	\mathrm{R}^n(GF)(X)/\mathrm{F}^1
	\simeq E_\infty^{0,n}
\]
tertanam dalam $E_2^{0,n}=G(\mathrm{R}^nF(X))$; komposisinya memberikan
morfisme tepi kedua. Arah ini sesuai dengan
\eqref{eqn:edge-ss-coh}.

Perhitungan dalam metode kerja di atas berlaku pada setiap barisan spektral
kuadran pertama dan memberi
\[
	0\to E_2^{1,0}\to H^1\to E_2^{0,1}
	\xrightarrow{d_2}E_2^{2,0}\to H^2.
\]
Substitusi
\[
	E_2^{p,q}=\mathrm{R}^pG(\mathrm{R}^qF(X)),
	\qquad H^n=\mathrm{R}^n(GF)(X)
\]
memberikan tepat barisan lima suku yang dinyatakan.

Jika $\mathrm{R}^qF(X)=0$ untuk $q>0$, hanya baris $q=0$ yang tak nol.
Tidak ada diferensial atau masalah ekstensi lintas baris, sehingga
\[
	E_2^{n,0}=E_\infty^{n,0}=\mathrm{R}^n(GF)(X),
\]
dan morfisme tepi pertama merupakan isomorfisme. Jika semua suku dengan $p>0$ lenyap,
hanya kolom $p=0$ yang tersisa; demikian pula
\[
	\mathrm{R}^n(GF)(X)=E_\infty^{0,n}=E_2^{0,n}
	=G(\mathrm{R}^nF(X)),
\]
dan morfisme tepi kedua merupakan isomorfisme.

% stable-exercise-id: o014.mastery.derived-spectral.002.ex07
\subsubsection*{Latihan 7 --- Kohomologi grup siklik secara konkret}
\hypertarget{o014.mastery.derived-spectral.002.ex07}{ }
\label{mastery:derived-spectral-002-ex07}

Misalkan $C_m=\lrangle{\sigma}$ adalah grup siklik berorde $m$, dan
$A=\Z/r\Z$ dilengkapi dengan aksi trivial dari $C_m$. Tuliskan $d=\gcd(m,r)$.
Gunakan resolusi periodik
\[
	\cdots\xrightarrow{\nu}\Z[C_m]
	\xrightarrow{\sigma-1}\Z[C_m]
	\xrightarrow{\nu}\Z[C_m]
	\xrightarrow{\sigma-1}\Z[C_m]\to\Z\to0,
	\qquad \nu=1+\sigma+\cdots+\sigma^{m-1},
\]
untuk menghitung $\Hm^n(C_m,A)$ untuk semua $n\geq0$. Sebagai pengecekan,
tentukan pula hasil untuk koefisien trivial $A=\Z$.

% stable-solution-id: o014.mastery.derived-spectral.002.sol07
\paragraph{Solusi.}
\hypertarget{o014.mastery.derived-spectral.002.sol07}{ }
Setiap $\Hom_{\Z[C_m]}(\Z[C_m],A)$ isomorfik secara kanonik dengan $A$.
Karena aksinya trivial, $\sigma-1$ bekerja sebagai nol, sedangkan $\nu$ bekerja
sebagai perkalian dengan $m$. Setelah menerapkan $\Hom_{\Z[C_m]}(-,A)$,
diperoleh kompleks
\[
	0\to\underbracket{A}_{0}\xrightarrow{0}
	\underbracket{A}_{1}\xrightarrow{\times m}A
	\xrightarrow{0}A\xrightarrow{\times m}A\to\cdots.
\]
Oleh karena itu,
\[
	\Hm^n(C_m,A)\simeq
	\begin{cases}
		A,&n=0,\\
		A[m]:=\Ker(\times m:A\to A),&n>0\text{ ganjil},\\
		A/mA,&n>0\text{ genap}.
	\end{cases}
\]
Dalam $A=\Z/r\Z$, subgrup $A[m]$ dibangkitkan oleh kelas $r/d$ dan
berorde $d$, sehingga $A[m]\simeq\Z/d\Z$. Subgrup $mA$ berorde $r/d$,
sehingga hasil bagi $A/mA$ juga siklik berorde $d$. Jadi
\[
	\Hm^n(C_m,\Z/r\Z)\simeq
	\begin{cases}
		\Z/r\Z,&n=0,\\
		\Z/d\Z,&n>0.
	\end{cases}
\]
Untuk $A=\Z$, kernel perkalian dengan $m$ adalah nol dan kokernelnya
adalah $\Z/m\Z$. Maka
\[
	\Hm^n(C_m,\Z)\simeq
	\begin{cases}
		\Z,&n=0,\\
		0,&n>0\text{ ganjil},\\
		\Z/m\Z,&n>0\text{ genap}.
	\end{cases}
\]
Hasil ini merupakan spesialisasi langsung dari
Proposisi~\ref{prop:group-cohomology-cyclic}.

% stable-exercise-id: o014.mastery.derived-spectral.002.ex08
\subsubsection*{Latihan 8 --- Pemeriksaan filtrasi multiplikatif}
\hypertarget{o014.mastery.derived-spectral.002.ex08}{ }
\label{mastery:derived-spectral-002-ex08}

Misalkan $(A^\bullet,d,\mu)$ aljabar-dg kohomologis di atas gelanggang
komutatif $\Bbbk$, dilengkapi filtrasi menurun yang menyeluruh dan terbatas
pada setiap derajat,
\[
	d(\mathrm{F}^pA)\subset\mathrm{F}^pA,
	\qquad
	\mathrm{F}^pA\cdot\mathrm{F}^{p'}A
	\subset\mathrm{F}^{p+p'}A.
\]
\begin{enumerate}[(i)]
	\item Buktikan bahwa $E_0=\gr_{\mathrm F}A$ mempunyai perkalian
bergradasi dan bahwa $d_0$ memenuhi aturan Leibniz.
	\item Jika $x$ dan $y$ mewakili kelas pada
$E_r^{p,q}$ dan $E_r^{p',q'}$, buktikan bahwa $xy$ mewakili kelas pada
$E_r^{p+p',q+q'}$, hasil perkaliannya tidak bergantung pada pilihan wakil, dan
	\[
		d_r(xy)=d_r(x)y+(-1)^{p+q}x\,d_r(y).
	\]
	\item Untuk filtrasi terinduksi
	$\mathrm{F}^p\Hm^n(A)=\Image[\Hm^n(\mathrm{F}^pA)\to\Hm^n(A)]$,
	buktikan
	\[
		\mathrm{F}^p\Hm^n(A)\cdot\mathrm{F}^{p'}\Hm^{n'}(A)
		\subset\mathrm{F}^{p+p'}\Hm^{n+n'}(A)
	\]
	dan bahwa isomorfisme konvergensi
	$E_\infty\simeq\gr_{\mathrm F}\Hm(A)$ mempertahankan perkalian.
\end{enumerate}

% stable-solution-id: o014.mastery.derived-spectral.002.sol08
\paragraph{Solusi.}
\hypertarget{o014.mastery.derived-spectral.002.sol08}{ }
Untuk kelas homogen
$[x]\in\gr^p_{\mathrm F}A^n$ dan
$[y]\in\gr^{p'}_{\mathrm F}A^{n'}$, definisikan
\[
	[x][y]:=[xy]\in\gr^{p+p'}_{\mathrm F}A^{n+n'}.
\]
Jika $x$ atau $y$ diubah oleh unsur pada tingkat filtrasi berikutnya,
kompatibilitas perkalian menempatkan perubahan pada hasil kali $xy$ dalam
$\mathrm{F}^{p+p'+1}$; jadi perkalian terdefinisi dengan baik. Karena $d$
mempertahankan filtrasi, ia menginduksi $d_0$, dan aturan Leibniz pada $A$
memberi
\[
	d_0([x][y])=d_0([x])\,[y]+(-1)^n[x]\,d_0([y]).
\]

Pada halaman $E_r$, wakil dapat dipilih sehingga
\[
	x\in\mathrm{F}^pA^{p+q},\quad
	dx\in\mathrm{F}^{p+r}A^{p+q+1},
\]
dan demikian pula
$y\in\mathrm{F}^{p'}A^{p'+q'}$ dengan
$dy\in\mathrm{F}^{p'+r}A^{p'+q'+1}$. Kompatibilitas filtrasi dan Leibniz
memberi
\[
	xy\in\mathrm{F}^{p+p'}A^{p+q+p'+q'},
\]
serta
\[
	d(xy)=dx\,y+(-1)^{p+q}x\,dy
	\in\mathrm{F}^{p+p'+r}A^{p+q+p'+q'+1}.
\]
Jadi $xy$ memang menentukan kelas pada $E_r^{p+p',q+q'}$, dan rumus
Leibniz bagi $d_r$ mengikuti langsung.

Untuk memeriksa bahwa hasilnya tidak bergantung pada pilihan wakil, cukup periksa perubahan oleh batas-$r$ dan
oleh unsur pada filtrasi lebih dalam. Jika, misalnya, $x$ diganti dengan
$x+du$, dengan
$u\in\mathrm{F}^{p-r+1}A^{p+q-1}$, maka
\[
	(du)y=d(uy)-(-1)^{p+q-1}u\,dy.
\]
Suku pertama merupakan batas, sedangkan suku kedua terletak dalam
\[
	\mathrm{F}^{p-r+1+p'+r}A
	=\mathrm{F}^{p+p'+1}A,
\]
sehingga tidak mengubah kelas yang bersangkutan. Perubahan wakil $y$
ditangani secara simetris; perubahan pada filtrasi lebih dalam bahkan lebih
langsung. Maka perkalian terdefinisi pada setiap halaman dan diwariskan ke
$E_{r+1}\simeq\Hm(E_r,d_r)$.

Terakhir, jika $[x]\in\mathrm{F}^p\Hm^n(A)$ dan
$[y]\in\mathrm{F}^{p'}\Hm^{n'}(A)$, pilih wakil tertutup
$x\in\mathrm{F}^pA^n$ dan $y\in\mathrm{F}^{p'}A^{n'}$. Produk $xy$ tertutup
dan terletak dalam $\mathrm{F}^{p+p'}A^{n+n'}$, sehingga inklusi filtrasi
yang diminta berlaku. Isomorfisme
$E_\infty^{p,q}\simeq\gr^p_{\mathrm F}\Hm^{p+q}(A)$ dibangun dari
wakil terfiltrasi yang sama pada kedua ruas. Hasil kali dua wakil dipetakan ke
kelas hasil kalinya; karena itu isomorfisme ini mempertahankan
perkalian. Pemeriksaan ini mewakili argumen dalam
Proposisi~\ref{prop:dga-mult-ss}.
